% Generated draft: requires mathematical and source review.
\documentclass[11pt]{article}
\usepackage[margin=25mm]{geometry}
\usepackage{fontspec}
\setmainfont{Latin Modern Roman}
\usepackage{amsmath,amssymb,mathtools,unicode-math}
\setmathfont{Latin Modern Math}
\usepackage{enumitem,needspace,xurl,hyperref}
\setlength{\parindent}{0pt}
\setlength{\parskip}{5pt}
\setlength{\emergencystretch}{3em}
\newcommand{\N}{\mathbb{N}}
\newcommand{\Z}{\mathbb{Z}}
\newcommand{\Q}{\mathbb{Q}}
\newcommand{\R}{\mathbb{R}}
\newcommand{\C}{\mathbb{C}}
\newcommand{\F}{\mathbb{F}}
\newcommand{\A}{\mathbb{A}}
\newcommand{\PP}{\mathbb{P}}
\newcommand{\E}{\mathbb{E}}
\newcommand{\eps}{\varepsilon}
\newcommand{\dd}{\,\mathrm{d}}
\begin{document}
\section*{UnsolvedMath ID: 30000266}
\textbf{Source problem code:} catalogue entry
\subsection{Definitions and mathematical statement}
For ellipticity constants $0<\lambda\le\Lambda$ and a real symmetric $N\times N$ matrix $X$ with eigenvalues $e_1,\ldots,e_N$, Pucci's maximal operator is
\[
 \mathcal M^+_{\lambda,\Lambda}(X)
   =\Lambda\sum_{e_j>0}e_j+\lambda\sum_{e_j<0}e_j.
\]
Let $p^*>1$ denote the critical exponent for this operator in the parameter regime and definition of the report. It is not silently identified with the Laplacian's Sobolev exponent.

The verified source asks two questions:
\begin{enumerate}
\item If $u\in C^2(\R^N)$, $u\ge0$, and
\[
 \mathcal M^+_{\lambda,\Lambda}(D^2u)+u^{p^*}=0\quad\text{on }\R^N,
\]
must $u$ be radial about some point? Explicitly, must there exist $x_0\in\R^N$ and a function $U:[0,\infty)\to[0,\infty)$ such that $u(x)=U(|x-x_0|)$?
\item For every $1<p<p^*$, must every nonnegative entire solution of
\[
 \mathcal M^+_{\lambda,\Lambda}(D^2u)+u^p=0\quad\text{on }\R^N
\]
be identically zero?
\end{enumerate}
The second statement is interpreted in the same classical $C^2$ class unless the report specifies another solution notion. The center in the first is not fixed to the origin, because translations preserve the equation. No boundedness, decay, finite-energy or boundary hypothesis is added.

The extract does not give the defining characterization of $p^*$, its precise dependence on $N,\lambda,\Lambda$, or the allowed dimension/ellipticity regime ensuring a finite critical exponent. Those are necessary source qualifications and cannot be filled by the familiar semilinear Laplacian formula.

The literature assessment checked on 21 August 2026 says its earlier extraction contained only ``Open Problem'' and classifies it conservatively. The corrected statement and supplied source verification now restore both actual questions. The earlier extraction-based assessment therefore does not verify their present status or provide a theorem resolving either one.
\subsection{Short English statement}
For Pucci's maximal operator, are all nonnegative critical entire solutions radial up to translation, and are all nonnegative subcritical entire solutions zero? The operator-specific critical exponent and its admissible parameter regime must come from the report.
\subsection{Sources}
\begin{itemize}
\item \emph{Topological and Variational Methods for Differential Equations} (2006), Oberwolfach Reports, DOI 10.4171/owr/2005/29, pp. 1601--1678, \url{https://doi.org/10.4171/owr/2005/29}; proposed year supplied as 2005.
\item The sole source in the assessment checked 21 August 2026 is the ``original problem report,'' \url{https://doi.org/10.4171/owr/2005/29}. That assessment refers to an earlier incomplete extraction rather than a resolved theorem about the restored equations.
\item Local UnsolvedMath registry, permanent ID $30000266$, exact source code \texttt{catalogue entry}, created 25 August 2026 and updated 16 September 2026. Supplied corrected-and-verified status confirms the maximal operator, critical-exponent role, radiality and subcritical Liouville targets.
\end{itemize}
\end{document}
